\chapter{Abstract}

Large-scale industrial software systems face significant challenges in executing regression test suites efficiently. Bach et al.\ report that in a commercial SAP environment, over 180,000 automated tests required approximately 35 hours to complete, creating substantial execution bottlenecks and infrastructure costs~\cite{bach2017}. Further analysis revealed that more than 80\,\% of these tests were redundant in terms of coverage---their coverage vectors were fully subsumed by other tests in the suite. This degree of redundancy leads to severe inefficiencies, wasting computational resources and delaying feedback in continuous-integration workflows. Test-suite minimisation addresses this challenge by systematically removing tests that are redundant with respect to defined adequacy criteria---such as statement coverage, branch coverage, or mutation score---while preserving fault-detection effectiveness~\cite{yoo2012,elbaum2001}. Formally, a test $t_i$ is considered redundant if its coverage vector $C(t_i)$ is subsumed by the union of coverage vectors from other selected tests: $C(t_i) \subseteq \bigcup_{t_j \in S} C(t_j)$~\cite{harrold1993}. However, purely coverage-driven reduction approaches risk excluding tests with unique fault-detection capabilities that are not captured by structural coverage alone, thereby compromising defect-revealing power~\cite{wong1995,rothermel1998}.

This thesis investigates test-suite minimisation within the \emph{Large-Scale Software Observatorium} (LASSO)~\cite{Kessel2022}, a research platform for automated selection, execution, and comparative analysis of software systems. LASSO integrates open-source repositories and combines static (code-level) and dynamic (execution-level) analyses within a unified framework designed for language independence. Its current implementation focuses on Java systems and enables empirically grounded studies of \emph{Big Code} ecosystems---large, diverse corpora of source code enriched with metadata such as bug-fix histories, version control traces, and execution outcomes, enabling statistical and probabilistic modeling techniques for software engineering research~\cite{allamanis2018}..

Our approach leverages \emph{Stimulus--Response Matrices} (SRMs)~\cite{Kessel2022}, a two-dimensional representation mapping test stimuli (rows) to observed execution outcomes (column entries). We extend SRM entries to incorporate coverage and mutation results, enabling minimization algorithms to operate on behavioral data rather than language-specific code structures. While current tooling supports Java systems, this abstraction is conceptually cross-language compatible.

Through a targeted literature survey encompassing greedy heuristics, exact optimization, search-based, and data-driven minimization strategies, we comparatively assessed candidate algorithms along six established criteria: coverage preservation, reduction effectiveness, fault-detection retention~\cite{yoo2012,zhang2010}, computational scalability, SRM compatibility, and implementation feasibility. The Mutation-Enhanced and Overlap-Aware Harrold--Gupta--Soffa (ME-OA-HGS) algorithm emerged as the most promising candidate. ME-OA-HGS extends the classical HGS heuristic~\cite{harrold1993} with mutation weighting and overlap penalties to balance coverage preservation and fault-detection power~\cite{bach2017}.

We successfully integrated ME-OA-HGS into LASSO and validated the implementation on real-world SRMs. Prior empirical studies report 45--70\,\% test suite reduction while maintaining $\geq$\,95\,\% statement coverage and $\geq$\,90\,\% mutation-score retention~\cite{harrold1993,bach2017}. The algorithm executes in polynomial time, ensuring computational feasibility for large-scale software analysis workflows. % TODO: Verify final complexity claim O(nm log m) vs O(n(n+m)·MAX_CARD) based on implementation chapter analysis before submission.